\documentclass[a4paper, 12pt]{article}
%\usepackage[skip]{parskip}
\usepackage{fullpage}
\usepackage{amsmath}
\usepackage{rotating}
\usepackage[IEEE]{linearcoefficients}

\begin{document}
\title{linearcoefficients -- version 1.0.0}
\author{Einar Halvorsen}
\date{May 25, 2026}
\maketitle

\abstract{
This package helps dealing with the notation for linear elastic/piezoelectric materials by defining commands for fields and coefficients as well as support for easy switching between different notations. There is also support for handling the superscripts denoting the conditions for the coefficients. 
}

\section{Notation}
This package provides commands for the symbols of the physical quantities of a piezoelectric as well as for the linear coefficients relating them. The actual symbols are defined according to specified a notation. The primary objective of this functionality is to add a layer of abstraction that make reuse of \LaTeX{} code in a context that requires different notation.

Some notations are predefined: One mostly congruent with the definitions in the IEEE standard on piezoelectricity \cite{IEEEpiezostd1996}. It is set with the command \verb|\IEEENotation| or by passing the option \verb|IEEE| to the package. The other notations follow Nye \cite{Nye1985}, Ikeda \cite{Ikeda1996} and Mason \cite{Mason1958}, The correspoding options are \verb|Nye|, \verb|Ikeda| and \verb|Mason| with respective commands \verb|\NyeNotation|, \verb|\IkedaNotation|, \verb|\MasonNotation|. 

Any number of notational options can be given to the package. The last one in the options list is the one chosen. If none is given, the package defaults to the IEEE-like notation. This defines the notation globally in the document.
The predefined notations are given in table \ref{tab:notations} which shows the physical quantity, its corresponding symbol in each predefined notation and the command to generate it.

The definitions can be changed within the local scope (group). For example, in a document where the globally defined notation is the IEEE-like notation, the code
\begin{verbatim}
\[
  {\NyeNotation\stress_{ik} = \stiffness_{ijkl}\strain_{kl}}
  \qquad\mbox{vs.}\qquad
  {\noelectrical \stress_{ik} = \stiffness_{ijkl}\strain_{kl}}.
\]
\end{verbatim}
produces
\[
  {\NyeNotation\stress_{ik} = \stiffness_{ijkl}\strain_{kl}}
  \qquad\mbox{vs.}\qquad
  {\noelectrical \stress_{ik} = \stiffness_{ijkl}\strain_{kl}}.
\]
\begin{table}
  \centering
  \caption{Symbols for  the predefined notations (sans indices). \label{tab:notations}}
  \begin{tabular}{llcccc}
    \hline\hline
    Quantity &  Command & IEEE & Nye & Ikeda & Mason\\
    \hline 
    electric field & \verb|\efield|
                        & $\IEEENotation\efield$
                               & $\NyeNotation\efield$
                                     & $\IkedaNotation\efield$
                                             & $\MasonNotation\efield$\\
    electric polarization & \verb|\polarization|
                        & $\IEEENotation\polarization$
                               & $\NyeNotation\polarization$
                                     & $\IkedaNotation\polarization$
                                             & $\MasonNotation\polarization$\\
    electric flux density & \verb|\dfield|
                        & $\IEEENotation\dfield$
                               & $\NyeNotation\dfield$
                                     & $\IkedaNotation\dfield$
                                             & $\MasonNotation\dfield$\\
    stress & \verb|\stress|
                        & $\IEEENotation\stress$
                               & $\NyeNotation\stress$
                                     & $\IkedaNotation\stress$
                                             & $\MasonNotation\stress$\\
    strain & \verb|\strain|
                        & $\IEEENotation\strain$
                               & $\NyeNotation\strain$
                                     & $\IkedaNotation\strain$
                                             & $\MasonNotation\strain$\\
    temperature & \verb|\temperature|
                        & $\IEEENotation\temperature$
                               & $\NyeNotation\temperature$
                                     & $\IkedaNotation\temperature$
                                             & $\MasonNotation\temperature$\\
    entropy density & \verb|\entropydensity|
                        & $\IEEENotation\entropydensity$
                               & $\NyeNotation\entropydensity$
                                     & $\IkedaNotation\entropydensity$
                                             & $\MasonNotation\entropydensity$\\
    \hline
    stiffness coefficient & \verb|\stiffness|
                        & $\IEEENotation\stiffness$
                               & $\NyeNotation\stiffness$
                                     & $\IkedaNotation\stiffness$
                                             & $\MasonNotation\stiffness$\\
    compliance coefficient & \verb|\compliance|
                        & $\IEEENotation\compliance$
                               & $\NyeNotation\compliance$
                                     & $\IkedaNotation\compliance$
                                             & $\MasonNotation\compliance$\\
    permittivity & \verb|\permittivity|
                        & $\IEEENotation\permittivity$
                               & $\NyeNotation\permittivity$
                                     & $\IkedaNotation\permittivity$
                                             & $\MasonNotation\permittivity$\\
    impermittivity & \verb|\impermittivity|
                        & $\IEEENotation\impermittivity$
                               & $\NyeNotation\impermittivity$
                                     & $\IkedaNotation\impermittivity$
                                             & $\MasonNotation\impermittivity$\\
    piezoelectric coefficient &\verb|\piezoh|
                        & $\IEEENotation\piezoh$
                               & $\NyeNotation\piezoh$
                                     & $\IkedaNotation\piezoh$
                                             & $\MasonNotation\piezoh$\\
    piezoelectric coefficient & \verb|\piezog|
                        & $\IEEENotation\piezog$
                               & $\NyeNotation\piezog$
                                     & $\IkedaNotation\piezog$
                                             & $\MasonNotation\piezog$\\
    piezoelectric coefficient & \verb|\piezoe|
                        & $\IEEENotation\piezoe$
                               & $\NyeNotation\piezoe$
                                     & $\IkedaNotation\piezoe$
                                             & $\MasonNotation\piezoe$\\
    piezoelectric coefficient& \verb|\piezod|
                        & $\IEEENotation\piezod$
                               & $\NyeNotation\piezod$
                                     & $\IkedaNotation\piezod$
                                             & $\MasonNotation\piezod$\\
    heat capacity & \verb|\heatcapacity|
                        & $\IEEENotation\heatcapacity$
                               &  $\NyeNotation\heatcapacity$
                                     &  $\IkedaNotation\heatcapacity$
                                             &  $\MasonNotation\heatcapacity$\\
    pyroelectric coefficient & \verb|\pyroelec|
                        & $\IEEENotation\pyroelec$
                               &  $\NyeNotation\pyroelec$
                                     &  $\IkedaNotation\pyroelec$
                                             &  $\MasonNotation\pyroelec$\\
    heat of polarization & \verb|\heatofpolarization|
                        & $\IEEENotation\heatofpolarization$
                               &  $\NyeNotation\heatofpolarization$
                                     &  $\IkedaNotation\heatofpolarization$
                                             & $\MasonNotation
                                               \heatofpolarization$\\
    coefficient of thermal expansion  & \verb|\thermalexpansion|
                        & $\IEEENotation\thermalexpansion$
                               &  $\NyeNotation\thermalexpansion$
                                     &  $\IkedaNotation\thermalexpansion$
                                             &  $\MasonNotation
                                               \thermalexpansion$\\
    heat of deformation  & \verb|\heatofdeformation|
                        & $\IEEENotation\heatofdeformation$
                               &  $\NyeNotation\heatofdeformation$
                                     &  $\IkedaNotation\heatofdeformation$
                                             &  $\MasonNotation\heatofdeformation$\\
    \hline
    mass density & \verb|\massdensity|
                        & $\IEEENotation\massdensity$
                               &  $\NyeNotation\massdensity$
                                     &  $\IkedaNotation\massdensity$
                                             &  $\MasonNotation\massdensity$\\
    \hline
  \end{tabular}
\end{table}


\section{Free  variables}
It is customary to denote the conditions under which a linear coefficient is measured by superscripts. For example, a stiffness constant measured at constant electric field and temperature is written ${\freeefield\freetemperature\stiffness_{ij}}$ in the IEEE-like notation. It is possible to accomplish this by writing \verb|\stiffness^{\efield,\temperature}_{ij}|, but the package provides notation to automate the inclusion of such superscripts. 

The automated notation hinges on declaration of at most one field variable from
each of the domains mechanical, electrical and thermal as free variables.
The package then includes the relevant variables among those that are declared as independent in the superscripts.

The command  \verb|\freestress| declares stress as a free variable in
the mechanical domain. Since the conjugate variable strain then must be dependent,
any prior declaration of it as free is nulled by the command.
Likewise \verb|\freestrain| declares strain as the free variable in the
mechanical domain. In the electrical domain the similar commands are 
\verb|\freeefield| and  \verb|\freedfield| for respectively the electrical field and the electric flux density. In the thermal domain the commands are
\verb|\freetemperature| and \verb|\freeentropydensity| for respectively temperature and entropy density.

The commands \verb|\nomechanical|, \verb|\noelectrical| and \verb|\nothermal| clear all declarations of free variables in the domain indicated by the name of the command. The command \verb|\nofree| clears all declarations of
free variables. Common combinations of free variables have macros as given in table \ref{tab:formcommands}.
\begin{table}
  \caption{Macros for common combinations of free variables and their equivalents in terms of simpler commands}
  \label{tab:formcommands}
  \centering
  \begin{tabular}{ll}
    \hline\hline
    \rule{0pt}{2.5ex}Macro  & Equivalent  \\
    \hline
    \verb|\hform| & \verb|\freestrain\freedfield\nothermal| \\          
    \verb|\gform| & \verb|\freestress\freedfield\nothermal| \\          
    \verb|\eform| & \verb|\freestrain\freeefield\nothermal| \\          
    \verb|\dform| & \verb|\freestress\freeefield\nothermal| \\          
    \verb|\Uform| & \verb|\freestrain\freedfield\freeentropydensity| \\ 
    \verb|\Hform| & \verb|\freestress\freeefield\freeentropydensity| \\ 
    \verb|\HMform| & \verb|\freestress\freedfield\freeentropydensity| \\
    \verb|\HEform| & \verb|\freestrain\freeefield\freeentropydensity| \\
    \verb|\Aform| & \verb|\freestrain\freedfield\freetemperature| \\    
    \verb|\Gform| & \verb|\freestress\freeefield\freetemperature| \\    
    \verb|\GMform| & \verb|\freestress\freedfield\freetemperature| \\   
    \verb|\GEform| & \verb|\freestrain\freeefield\freetemperature| \\   
    \hline
  \end{tabular}
\end{table}

\begin{table}
  \caption{Examples of IEEE-like notation for coefficients depending on which fields are declared free.}
  \label{tab:coefficicents}
  \newcommand{\coefficientlist}{$\stiffness_{ijkl}$, $\compliance_{ijkl}$,
         $\permittivity_{ij}$, $\impermittivity_{ij}$,
         $\piezoh_{ikl}$,  $\piezog_{ikl}$, $\piezoe_{ikl}$, $\piezod_{ikl}$}
  \centering
  \begin{tabular}{lc}
    \hline\hline
    Free & Coefficients\\
    \hline
    \rule{0pt}{2.5ex} -- & \coefficientlist\\
    \rule{0pt}{2.5ex} $\efield$ & \freeefield\coefficientlist\\
    \rule{0pt}{2.5ex} $\dfield$ & \freedfield\coefficientlist\\
    \rule{0pt}{2.5ex} $\efield$, $\stress$
         & \freeefield\freestress\coefficientlist\\
    \rule{0pt}{2.5ex} $\efield$, $\strain$
         & \freeefield\freestrain\coefficientlist\\
    \rule{0pt}{2.5ex} $\efield$, $\strain$, $\temperature$
         & \freeefield\freestrain\freetemperature\coefficientlist\\
    \rule{0pt}{2.5ex} $\efield$, $\strain$, $\entropydensity$
         & \freeefield\freestrain\freeentropydensity\coefficientlist\\
    \hline
  \end{tabular}
\end{table}

As an example usage, consider the code
\begin{verbatim}
\freestress\freeefield\freetemperature
{\NyeNotation
\begin{align*}
  \strain_{ij} & = \compliance_{ijkl}\stress_{kl} +
                 \piezod_{kij}\efield_k
                 + \thermalexpansion_{ij}\temperature \\
  \dfield_i & = \piezod_{ikl}\stress_{kl}
              + \permittivity_{ik}\efield_k
              + \pyroelec_i\temperature \\
  \Delta\entropydensity & = \thermalexpansion_{kl}\stress_{kl}
                          + \pyroelec_k\efield_k
                          + (\heatcapacity/\temperature)
                          \Delta\temperature.
\end{align*}
}
\end{verbatim}
for a set of linear constitutive equations. It renders as
\freestress\freeefield\freetemperature
{\NyeNotation
\begin{align*}
  \strain_{ij} & = \compliance_{ijkl}\stress_{kl} + \piezod_{kij}\efield_k + \thermalexpansion_{ij}\Delta\temperature \\
  \dfield_i & = \piezod_{ikl}\stress_{kl} + \permittivity_{ik}\efield_k + \pyroelec_i\Delta\temperature \\
  \Delta\entropydensity & = \thermalexpansion_{kl}\stress_{kl} + \pyroelec_k\efield_k + (\heatcapacity/\temperature)\Delta\temperature.
\end{align*}
}

\section{Matrix notation}
There is support for matrix notation where rectangular matrices are indicated
by double underscores and column matrices are indicated by single underscores
provided respectively by the commands \verb|\mat| and \verb|\colmat|. There are also varieties of these called \verb|\matp| and \verb|\colmatp| that put a prime on the symbol in questions. For example
\begin{verbatim}
\IkedaNotation\nofree
\begin{align*}
  \colmat{\dfield} = \mat{\permittivity}\colmat{\efield},
  && \colmatp{\dfield} = \mat{a}\colmat{\dfield},
  && \colmatp{\efield} = \mat{a}\colmat{\efield},   
  && \matp{\permittivity} = \mat{a}^\mathrm{T}\mat{\permittivity}\mat{a}.    
\end{align*}
\end{verbatim}
generates
{\IkedaNotation\nofree
\begin{align*}
  \colmat{\dfield} = \mat{\permittivity}\colmat{\efield},
  && \colmatp{\dfield} = \mat{a}\colmat{\dfield},
  && \colmatp{\efield} = \mat{a}\colmat{\efield},   
  && \matp{\permittivity} = \mat{a}^\mathrm{T}\mat{\permittivity}\mat{a}.    
\end{align*}
}

The commad \verb|\umat| automatically generates correct underlines on the coefficient or field quantity but avoids underlining superscripts. Examples in Ikeda's notation are
{\IkedaNotation
\begin{align}
  \umat\polarization,\quad && && \umat\efield,\quad && \umat\dfield,\quad
  &&  \umat\stress,\quad && \umat\strain,\quad
  && \umat\temperature,\quad &&  \umat\entropydensity,\nonumber \\
  {\nofree\umat\stiffness},\quad && {\GEform\umat\stiffness},\quad
  && {\nofree\umat\compliance},\quad && {\Gform\umat\compliance},\quad
  && {\nofree\umat\permittivity},\quad && {\GEform\umat\permittivity},\quad  
  && {\nofree\umat\impermittivity},\quad && {\GEform\umat\impermittivity},
  \nonumber \\
  {\nofree\umat\piezoh},\quad && {\GEform\umat\piezoh},\quad
  && {\nofree\umat\piezog},\quad && {\GEform\umat\piezog},\quad
  && {\nofree\umat\piezoe},\quad && {\GEform\umat\piezoe},\quad
  && {\nofree\umat\piezod},\quad && {\GEform\umat\piezod}, \nonumber \\
  {\nofree\umat\heatcapacity},\quad && {\GEform\umat\heatcapacity},\quad
  && {\nofree\umat\pyroelec},\quad && {\GEform\umat\pyroelec},\quad
  && {\nofree\umat\heatofpolarization},\quad && {\GEform\umat\heatofpolarization},\quad
  && {\nofree\umat\thermalexpansion},\quad && {\GEform\umat\thermalexpansion},\quad
  && {\nofree\umat\heatofdeformation},\quad && {\GEform\umat\heatofdeformation}.\quad\nonumber
\end{align}
}
The last line was generated by
\begin{verbatim}
  {\nofree\umat\heatcapacity},\quad
  && {\GEform\umat\heatcapacity},\quad
  && {\nofree\umat\pyroelec},\quad && {\GEform\umat\pyroelec},\quad
  && {\nofree\umat\heatofpolarization},\quad
  && {\GEform\umat\heatofpolarization},\quad
  && {\nofree\umat\thermalexpansion},\quad
  && {\GEform\umat\thermalexpansion},\quad
  && {\nofree\umat\heatofdeformation},\quad
  && {\GEform\umat\heatofdeformation}.
\end{verbatim}

\section{Defining own notation}
The user can define own notation by redefining the underlying symbols.
For example, if we want capital $C$ instead of lower case $c$ for stiffness defined in the standard notations,
we can write \verb|\renewcommand{\stiffnesssymbol}{C}|. This is how the standard notations are defined. As an example, the definition of \verb|\IkedaNotation| is
\begin{verbatim}
\newcommand{\IkedaNotation}{% (Ikeda's notation)
\renewcommand{\polarizationsymbol}{P}
\renewcommand{\efieldsymbol}{E}
\renewcommand{\dfieldsymbol}{D}
\renewcommand{\stresssymbol}{T}
\renewcommand{\strainsymbol}{S}
\renewcommand{\temperaturesymbol}{\Theta}
\renewcommand{\entropydensitysymbol}{\sigma}
\renewcommand{\stiffnesssymbol}{c}
\renewcommand{\compliancesymbol}{s}
\renewcommand{\permittivitysymbol}{\epsilon}
\renewcommand{\impermittivitysymbol}{\beta}
\renewcommand{\piezohsymbol}{h}
\renewcommand{\piezogsymbol}{g}
\renewcommand{\piezoesymbol}{e}
\renewcommand{\piezodsymbol}{d}
\renewcommand{\heatcapacitysymbol}{C}
\renewcommand{\pyroelecsymbol}{p}
\renewcommand{\heatofpolarizationsymbol}{q}
\renewcommand{\thermalexpansionsymbol}{\alpha}
\renewcommand{\heatofdeformationsymbol}{f}
\renewcommand{\massdensitysymbol}{\rho}
}
\end{verbatim}

\appendix
\section{Thermodynamic potentials and equations of state}
In this section, we rely on Mason's description in \cite{Mason1958} but use Ikeda's  notation. The constitutive equations can be derived from the thermodynamic potentials in table~\ref{tab:thermodyn_potentials}.

{\IkedaNotation
    {\clearpage\global\pdfpageattr\expandafter{\the\pdfpageattr/Rotate 90}
  \begin{sidewaystable}%[ht]
    \caption{Definitions of thermodynamic potential (densities) according to Mason \cite{Mason1958}.}
    \label{tab:thermodyn_potentials}
    \begin{tabular}{lclll}
      \hline\hline\rule{0pt}{2.5ex}Name
      & Symbol & Free variables & Definition & Differential \\
      \hline
      \rule{0pt}{2.5ex}Internal energy
           & $U$ & $\strain_{ij}$, $\dfield_k$, $\entropydensity$
                            & --
                            &  $\mathrm{d}U = \stress_{ij}\mathrm{d}\strain_{ij}
                              + \efield_k\mathrm{d}\dfield_k
                              + \temperature\mathrm{d}\entropydensity$\\
      Enthalpy
           & $H$ & $\stress_{ij}$, $\efield_k$, $\entropydensity$
                                & $U - \stress_{ij}\strain_{ij}
                                  - \efield_k\dfield_k$
                            & $\mathrm{d}H = -\strain_{ij}\mathrm{d}\stress_{ij}
                              - \dfield_k\mathrm{d}\efield_k
                              + \temperature\mathrm{d}\entropydensity$\\
      Elastic enthalpy
           & $H_1$ & $\stress_{ij}$, $\dfield_k$, $\entropydensity$
                                & $U - \stress_{ij}\strain_{ij}$
                            & $\mathrm{d}H_1 = -\strain_{ij}\mathrm{d}\stress_{ij}
                              + \efield_k\mathrm{d}\dfield_k
                              + \temperature\mathrm{d}\entropydensity$\\
      Electric enthalpy
           & $H_2$ & $\strain_{ij}$, $\efield_k$, $\entropydensity$
                                & $U - \efield_k\dfield_k$
                            & $\mathrm{d}H_2 = \stress_{ij}\mathrm{d}\strain_{ij}
                              - \dfield_k\mathrm{d}\efield_k
                              + \temperature\mathrm{d}\entropydensity$\\
      Free energy & $A$ & $\strain_{ij}$, $\dfield_k$, $\temperature$
                            & $U - \temperature\entropydensity$
                            &  $\mathrm{d}A = \stress_{ij}\mathrm{d}\strain_{ij}
                              + \efield_k\mathrm{d}\dfield_k
                              - \entropydensity\mathrm{d}\temperature$\\
      Gibbs function
           & $G$ & $\stress_{ij}$, $\efield_k$, $\temperature$
                                & $A - \stress_{ij}\strain_{ij}
                                  - \efield_k\dfield_k$
                            & $\mathrm{d}G = -\strain_{ij}\mathrm{d}\stress_{ij}
                              - \dfield_k\mathrm{d}\efield_k
                              -\entropydensity\mathrm{d}\temperature$\\
      Elastic Gibbs function
           & $G_1$ & $\stress_{ij}$, $\dfield_k$, $\temperature$
                                & $A - \stress_{ij}\strain_{ij}$
                            & $\mathrm{d}G_1 = -\strain_{ij}\mathrm{d}\stress_{ij}
                              + \efield_k\mathrm{d}\dfield_k
                              - \entropydensity\mathrm{d}\temperature$\\
      Electric Gibbs function
           & $G_2$ & $\strain_{ij}$, $\efield_k$, $\temperature$
                                & $A - \efield_k\dfield_k$
                            & $\mathrm{d}G_2 = \stress_{ij}\mathrm{d}\strain_{ij}
                              - \dfield_k\mathrm{d}\efield_k
                              - \entropydensity\mathrm{d}\temperature$\\
      \hline 
    \end{tabular}
  \end{sidewaystable}
  \clearpage\global\pdfpageattr\expandafter{\the\pdfpageattr/Rotate 0}}

From the table, we see that the fields can be given in terms of derivatives of the thermodynamic potentials as
\begin{align}
  \stress_{ij} & = \frac{\partial U}{\partial\strain_{ij}}
                 = \frac{\partial H_2}{\partial\strain_{ij}}
                 = \frac{\partial A}{\partial\strain_{ij}}
                 = \frac{\partial G_2}{\partial\strain_{ij}},\\
  \efield_{k} & = \frac{\partial U}{\partial\dfield_{k}}
                 = \frac{\partial H_1}{\partial\dfield_{k}}
                 = \frac{\partial A}{\partial\dfield_{k}}
                 = \frac{\partial G_1}{\partial\dfield_{k}},\\
  \temperature & = \frac{\partial U}{\partial\entropydensity}
                 = \frac{\partial H}{\partial\entropydensity}
                 = \frac{\partial H_1}{\partial\entropydensity}
                 = \frac{\partial H_2}{\partial\entropydensity},\\
  \strain_{ij} & = -\frac{\partial H}{\partial\stress_{ij}}
                 = -\frac{\partial H_1}{\partial\stress_{ij}}
                 = -\frac{\partial G}{\partial\stress_{ij}}
                 = -\frac{\partial G_1}{\partial\stress_{ij}},\\
  \dfield_{k} & = -\frac{\partial H}{\partial\efield_{k}}
                = -\frac{\partial H_2}{\partial\efield_{k}}
                = -\frac{\partial G}{\partial\efield_{k}}
                = -\frac{\partial G_2}{\partial\efield_{k}},\\
  \entropydensity & = -\frac{\partial A}{\partial\temperature}
                    = -\frac{\partial G}{\partial\temperature}
                    = -\frac{\partial G_1}{\partial\temperature}
                    = -\frac{\partial G_2}{\partial\temperature}.
\end{align}

The linear coefficients can then be defined as the first derivatives of a field under
certain conditions on the other free variables or as a second derivative of a potential.
For stiffness and compliance, we have
\begin{align}
  {\freedfield\freeentropydensity \stiffness_{ijkl}}
  & = \frac{\partial^2U}{\partial\strain_{kl}\partial\strain_{ij}},
  &
  {\freeefield\freeentropydensity \stiffness_{ijkl}}
  & = \frac{\partial^2H_2}{\partial\strain_{kl}\partial\strain_{ij}},
  &
  {\freedfield\freetemperature \stiffness_{ijkl}}
  & = \frac{\partial^2A}{\partial\strain_{kl}\partial\strain_{ij}},
  &
  {\freeefield\freetemperature \stiffness_{ijkl}}
  & = \frac{\partial^2G_2}{\partial\strain_{kl}\partial\strain_{ij}},\\
  {\freeefield\freeentropydensity \compliance_{ijkl}}
  & = -\frac{\partial^2H}{\partial\stress_{kl}\partial\stress_{ij}},
  &
  {\freedfield\freeentropydensity \compliance_{ijkl}}
  & = -\frac{\partial^2H_1}{\partial\stress_{kl}\partial\stress_{ij}},    
  &
  {\freeefield\freetemperature \compliance_{ijkl}}
  & = -\frac{\partial^2G}{\partial\stress_{kl}\partial\stress_{ij}},  
  &
  {\freedfield\freetemperature \compliance_{ijkl}}
  & = -\frac{\partial^2G_1}{\partial\stress_{kl}\partial\stress_{ij}}.    
\end{align}
Impermittivities and permittivities are
\begin{align}
  {\freestrain\freeentropydensity\impermittivity_{kl}}
  & = \frac{\partial U}{\partial\dfield_{k}\partial\dfield_{l}},
    &
  {\freestress\freeentropydensity\impermittivity_{kl}}
  & = \frac{\partial H_1}{\partial\dfield_{k}\partial\dfield_{l}},
    &
  {\freestrain\freetemperature\impermittivity_{kl}}
  & = \frac{\partial A}{\partial\dfield_{k}\partial\dfield_{l}},
    &
  {\freestress\freetemperature\impermittivity_{kl}}
  & = \frac{\partial G_1}{\partial\dfield_{k}\partial\dfield_{l}},
  \\
  {\freestress\freeentropydensity\permittivity_{kl}}
  & = -\frac{\partial^2 H}{\partial\efield_{k}\partial\efield_{l}},
  &  
    {\freestrain\freeentropydensity\permittivity_{kl}}
  & = -\frac{\partial^2 H_2}{\partial\efield_{k}\partial\efield_{l}},
  &
    {\freestress\freetemperature\permittivity_{kl}}
  & = -\frac{\partial^2 G}{\partial\efield_{k}\partial\efield_{l}},
  &
    {\freestrain\freetemperature\permittivity_{kl}}
  & = -\frac{\partial^2 G_2}{\partial\efield_{k}\partial\efield_{l}}.
\end{align}
We note that without specifying conditions, the specific heat capacity ${\nofree\heatcapacity}$ is given by
\begin{equation}
  {\nofree
  \frac{\massdensity\heatcapacity}{\temperature}
  = \frac{\mathrm{d}\entropydensity}{\mathrm{d}\temperature}}
\end{equation}
hence the different specific heat capacities can be determined as
\begin{align}
  {\freestrain\freedfield
  \frac{\massdensity\heatcapacity}{\temperature}}
  & = -\frac{\partial^2 A}{\partial\temperature^2},
  &
    {\freestress\freeefield
    \frac{\massdensity\heatcapacity}{\temperature}}
  & = -\frac{\partial^2 G}{\partial\temperature^2},
    &
      {\freestress\freedfield
      \frac{\massdensity\heatcapacity}{\temperature}}
  & = -\frac{\partial^2 G_1}{\partial\temperature^2},
    &
    {\freestrain\freeefield
      \frac{\massdensity\heatcapacity}{\temperature}}
  & = -\frac{\partial^2 G_2}{\partial\temperature^2}.
\end{align}
There is no common terminology for the reciprocal of the heat capacity,
so we do not consider the four other relations corresponding to that.

The piezoelectric coupling constants are
\begin{align}
  {\freeentropydensity\piezoh_{kij}}
  & = -\frac{\partial^2 U}{\partial\dfield_{k}\partial\strain_{ij}},
  &
  {\freetemperature\piezoh_{kij}}
  & = -\frac{\partial^2 A}{\partial\dfield_{k}\partial\strain_{ij}}, \\
  {\freeentropydensity\piezog_{kij}}
  & = -\frac{\partial^2 H_1}{\partial\dfield_{k}\partial\stress{ij}},
  &
   {\freetemperature\piezog_{kij}}  
  & = -\frac{\partial G_1}{\partial\dfield_{k}\partial\stress{ij}},
  \\
  {\freeentropydensity\piezoe_{kij}}
  & = -\frac{\partial^2 H_2}{\partial\efield_{k}\partial\strain_{ij}},
  &
    {\freetemperature\piezoe_{kij}}
  &  = -\frac{\partial^2 G_2}{\partial\efield_{k}\partial\strain_{ij}},  
  \\
  {\freeentropydensity\piezod_{kij}}
  & = -\frac{\partial H}{\partial\efield_{k}\partial\stress{ij}},
  &
    {\freetemperature\piezod_{kij}}
  & = -\frac{\partial G}{\partial\efield_{k}\partial\stress{ij}}.
\end{align}

The pyroelectric coefficients are
\begin{align}
  {\freestress\pyroelec_{k}}
  & =  -\frac{\partial^2 G}{\partial\efield_{k}\partial\temperature},
  &                
    {\freestrain\pyroelec_{k}}
  & =  -\frac{\partial^2 G_2}{\partial\efield_{k}\partial\temperature}.
\end{align}



}
\bibliographystyle{plain}
\bibliography{piezo, books}

\end{document}
